% answers_short.tex —— 2015 年·陕西卷·理工农医类 · 简短答案速查
% 由 tools/make_short_answers.py 自 answers.tex 的 \gksolhead/\gkanswer 提取，
% 并与 2015/2015_数学真题_陕西理_答案要点.md 逐题比对后生成，勿手改（改源 MD 后重跑）。

\begin{gknotice}
\textbf{使用说明}：本节只给最终答案便于快速核对，完整推导见解析版。答案为编写者独立推导，\textbf{非官方标准答案}。
\end{gknotice}

\gkpart{选择题}
\begin{center}\begin{tabular}{ccccccccccccc}
\toprule
题号 & 1 & 2 & 3 & 4 & 5 & 6 & 7 & 8 & 9 & 10 & 11 & 12 \\
\midrule
答案 & A & C & C & B & D & A & B & B & C & D & D & A \\
\bottomrule
\end{tabular}\end{center}

\gkpart{填空题}
\begin{itemize}[leftmargin=2.6em]
  \item[\textbf{13.}] 5。
  \item[\textbf{14.}] $2\sqrt2$。
  \item[\textbf{15.}] $(1,1)$。
  \item[\textbf{16.}] $6/5$（相同流速条件）。
\end{itemize}

\gkpart{解答题}
\begin{itemize}[leftmargin=2.6em]
  \item[\textbf{17.}] (1) $A=\pi/3$；(2) $3\sqrt3/2$
  \item[\textbf{18.}] (1) 证明见解析；(2) $\sqrt6/3$
  \item[\textbf{19.}] (1) 以样本频率估计分布列 $(0.2,0.3,0.4,0.1)$；$ET=32$ 分钟；(2) 未定；若往返独立同分布则为 $0.91$
  \item[\textbf{20.}] (1) $\sqrt3/2$；(2) $x^2/12+y^2/3=1$
  \item[\textbf{21.}] (1) 证明见解析；(2) $f_n\le g_n$，当 $x=1$ 取等
  \item[\textbf{22.}] (1) 证明见解析；(2) 直径为 $3$
  \item[\textbf{23.}] (1) $x^2+(y-\sqrt3)^2=3$；(2) $(3,0)$
  \item[\textbf{24.}] (1) $a=-3,b=1$；(2) 最大值 $4$
\end{itemize}
